% Generated draft: requires mathematical and source review.
\documentclass[11pt]{article}
\usepackage[margin=25mm]{geometry}
\usepackage{fontspec}
\setmainfont{Latin Modern Roman}
\usepackage{amsmath,amssymb,mathtools,unicode-math}
\setmathfont{Latin Modern Math}
\usepackage{enumitem,needspace,xurl,hyperref}
\setlength{\parindent}{0pt}
\setlength{\parskip}{5pt}
\setlength{\emergencystretch}{3em}
\newcommand{\N}{\mathbb{N}}
\newcommand{\Z}{\mathbb{Z}}
\newcommand{\Q}{\mathbb{Q}}
\newcommand{\R}{\mathbb{R}}
\newcommand{\C}{\mathbb{C}}
\newcommand{\F}{\mathbb{F}}
\newcommand{\A}{\mathbb{A}}
\newcommand{\PP}{\mathbb{P}}
\newcommand{\E}{\mathbb{E}}
\newcommand{\eps}{\varepsilon}
\newcommand{\dd}{\,\mathrm{d}}
\begin{document}
\section*{UnsolvedMath ID: 30002533}
\textbf{Source problem code:} catalogue entry
\subsection{Definitions and mathematical statement}
Let \(\mathbb H=\{z\in\C:\operatorname{Im}z>0\}\). A Hilbert modular surface is an arithmetic quotient of \(\mathbb H\times\mathbb H\) associated with a real quadratic field, with the two factors corresponding to its two real embeddings. The source does not specify the field, arithmetic subgroup, level structure, or treatment of quotient singularities and compactification.

A Kobayashi geodesic curve is understood here as a complex curve whose lifted parametrization from its universal covering disk is a holomorphic isometric embedding for the Kobayashi distances. The precise class of curves intended by the source, including any algebraicity or finite-volume requirement, remains unspecified.

A Lyapunov exponent describes an asymptotic exponential growth rate of vectors under a linear cocycle over a dynamical flow. Its definition requires the cocycle, flow, invariant measure, and time normalization. These data, and the particular exponent attached to a geodesic curve, are not identified in the supplied record. Write \(\lambda(C)\) provisionally for that intended exponent, without selecting these missing data.

\textbf{Source question, with unresolved qualifications.} Can one construct a Hilbert modular surface \(X\) and a Kobayashi geodesic curve \(C\subset X\) such that the intended Lyapunov exponent satisfies
\[
\lambda(C)\in\Q\setminus
\left\{1,\frac12,\frac13,\frac15,\frac17\right\}?
\]
This expresses the request for an exponent outside the displayed list; it does not assert realizability of every rational number outside that list. The source does not say whether the surface may vary freely, whether a particular surface is intended, or whether one distinguished exponent or any member of a Lyapunov spectrum is meant. Consequently, the displayed formulation is provisional until these qualifications and the defining dynamical data are recovered.
\subsection{Short English statement}
Construct a Kobayashi geodesic curve on a Hilbert modular surface with a rational Lyapunov exponent outside \(\{1,1/2,1/3,1/5,1/7\}\). The supplied record leaves the relevant dynamical system, exponent normalization, and admissible class of curves and surfaces unspecified.
\subsection{Sources}
\begin{itemize}
\item Local UnsolvedMath registry, permanent ID \(30002533\), source problem code \texttt{catalogue entry}, title \emph{New Rational Lyapunov Exponents on Hilbert Modular Surfaces}; collection: Oberwolfach Reports Open Problems. The supplied record assigns the problem to 2014 and names no individual proposer.
\item Supplied source citation: \emph{Mini-Workshop: Negative Curves on Algebraic Surfaces} (2014), Oberwolfach Reports, pp. 555--593, DOI 10.4171/owr/2014/10. The background describes these reports as documenting mathematical workshops and their open-problem discussions. \url{https://doi.org/10.4171/owr/2014/10}.
\item Supplied literature reference: M\"oller, 2012. \url{https://arxiv.org/abs/1111.2624}. No title or more specific bibliographic attribution is supplied.
\item Supplied literature reference: Bogomolov et al., 2025. \url{https://arxiv.org/abs/2512.09188}. No title or complete author list is supplied.
\item Literature assessment dated exactly \texttt{2026-08-21}: the supplied curation labels the problem open and states that known rational Lyapunov exponents for Kobayashi geodesics do not exhaust the possibilities, while construction of curves with further exponents remains unresolved. This is supplied dated curation derived from the OWR record, not an independently verified assessment. Its time-sensitive status and cited claims require independent verification.
\end{itemize}
\end{document}
